\subsection{A rigorous rational certificate for the new conjecture}

The floating-point checks above can be strengthened to an exact
residual enclosure without resolving the equality itself.

\begin{theorem}[Euler remainder for harmonic numerators]\label{thm:euler-certificate}
Let \(a,b,q\) be positive integers,
\(c_n=H_{qn}^{(b)}/(2n+1)^a\), and
\(L=\sum_{n\geq0}(-1)^nc_n\).
Define \(\Delta c_n=c_{n+1}-c_n\) and
\[
 E_N=\sum_{k=0}^{N-1}\frac{(-\Delta)^k c_0}{2^{k+1}},
 \qquad N\geq1.
\]
Then
\begin{equation}\label{eq:euler-harmonic-bound}
 0\leq E_N-L\leq (N+1)2^{-N}c_1
     =(N+1)2^{-N}\frac{H_q^{(b)}}{3^a}.
\end{equation}
\end{theorem}

\begin{proof}
On the unit square use the positive measure
\[
 \dd\mu(t,u)=
 \frac{(-\log t)^{a-1}(-\log u)^{b-1}}
 {(a-1)!(b-1)!(1-u)}\dd t\,\dd u.
\]
Set \(v=t^2\), \(z=t^2u^q\). Integral representations for reciprocal
integer powers give
\[
 c_n=\int(v^n-z^n)\dd\mu,\qquad
 \int(v-z)\dd\mu=c_1<\infty.
\]
The measure itself need not have finite mass. All differences below
are integrable because they are bounded by a constant times \(v-z\).
For \(k\geq1\),
\[
 (-\Delta)^k c_0
 =\int[(1-v)^k-(1-z)^k]\dd\mu\leq0,\qquad
 |(-\Delta)^k c_0|\leq k c_1.
\]
Thus the Euler series converges absolutely.
It equals \(L\): introduce an Abel factor into the original series,
sum the geometric differences under an integrable bound, and let
the factor increase to 1. The ordinary series converges absolutely
for \(a>1\). When \(a=1\), its magnitudes eventually decrease to zero:
the harmonic increment divided by \(H_{qn}^{(b)}\) is
\(O(1/(n\log n))\) for \(b=1\) and \(O(n^{-b})\) for \(b>1\),
whereas the denominator ratio is \(1+2/(2n+1)\).
Abel's theorem identifies the limits.
Every omitted Euler term is nonpositive, and
\[
 \sum_{k=N}^{\infty}\frac{k}{2^{k+1}}=(N+1)2^{-N}.
\]
This proves both the sign and the bound.
\end{proof}

The finite sum has the exact rational form
\begin{equation}\label{eq:finite-euler-integer}
 E_N=2^{-N}\sum_{n=0}^{N-1}(-1)^nc_n
             \sum_{j=n+1}^{N}\binom Nj.
\end{equation}
Indeed,
\(\sum_{k=n}^{N-1}\binom{k}{n}2^{-k-1}
=2^{-N}\sum_{j=n+1}^{N}\binom Nj\).
When \(q=1\) or \(2\), a common denominator is
\(2^N\operatorname{lcm}(1,2,\ldots,2N)^{a+b}\).
Thus all four nonclassical constants in the \(S_6\) basket have
rapid exact rational enclosures.

For \(c_n=(dn+1)^{-s}\), \(d=1\) or \(2\), a positive
moment measure of total mass 1 instead gives
\(0\leq L-E_N\leq2^{-N}\).
This encloses \(\eta(s)\), \(\beta(s)\), and hence
\(\zeta(3),\zeta(5),\log2\) using rational factors.
For \(\pi\) use Machin's identity
\[
 \pi=16\arctan(1/5)-4\arctan(1/239)
\]
and the alternating-series bound for each arctangent.

The delivered standard-library script \texttt{certify\_s6.py} uses
\(N=900\), 240 arctangent terms at \(1/5\), and 80 at \(1/239\).
It performs every interval operation with integers and rational
fractions, and rounds final decimal endpoints outward.
The resulting integer-normalized residual \(R_6\) satisfies the
deliberately widened, easily readable rational enclosure
\begin{equation}\label{eq:S6-certificate}
 -5.016\times10^{-261}\leq R_6\leq3.845\times10^{-262}.
\end{equation}
The full 280-place endpoints are recorded in the JSON certificate.
The exact arithmetic gate checks, in particular,
\[
 \boxed{|R_6|<10^{-260}.}
\]
This is a rigorous error statement about the residual.
It is still not a proof that the residual equals zero.

